\documentclass[11pt,a4paper]{amsart}
\usepackage[T1]{fontenc}
\usepackage{lmodern,amssymb,mathtools,mathrsfs,microtype,booktabs}
\usepackage[margin=28mm]{geometry}
\usepackage{needspace}
\PassOptionsToPackage{hyphens}{url}
\usepackage[hidelinks]{hyperref}
\urlstyle{same}
\numberwithin{equation}{section}
\newtheorem{theorem}{Theorem}[section]
\newtheorem{proposition}[theorem]{Proposition}
\newtheorem{lemma}[theorem]{Lemma}
\newtheorem{corollary}[theorem]{Corollary}
\newtheorem{conjecture}[theorem]{Conjecture}
\theoremstyle{definition}
\newtheorem{definition}[theorem]{Definition}
\theoremstyle{remark}
\newtheorem{remark}[theorem]{Remark}
\DeclareMathOperator{\tr}{tr}
\DeclareMathOperator{\diag}{diag}
\DeclareMathOperator{\Sym}{Sym}
\DeclareMathOperator{\Scal}{Scal}
\DeclareMathOperator{\Vol}{Vol}
\DeclareMathOperator{\diver}{div}
\newcommand{\R}{\mathbb R}
\newcommand{\Sph}{\mathbb S}
\newcommand{\K}{\mathcal K}
\newcommand{\T}{\mathsf T}
\newcommand{\N}{\mathsf N}
\newcommand{\Pz}{\mathsf P_0}
\newcommand{\av}[1]{\langle #1\rangle_M}
\newcommand{\Szero}{\operatorname{Sym}^2_0\R^4}
\newcommand{\Hfour}{\mathscr H_4}
\setlength{\emergencystretch}{1.5em}
\allowdisplaybreaks[2]
\title[The second pinching problem in dimension four]{A proof of the second pinching theorem\\
in dimension four}
\author{Jinxi Zhang}
\address{University of Cambridge, Cambridge, United Kingdom}
\date{18 October 2026}
\subjclass[2020]{53C42, 53C24}
\keywords{Minimal hypersurface, constant scalar curvature, Chern conjecture, second pinching problem}
\hypersetup{pdftitle={A proof of the second pinching theorem in dimension four},
 pdfsubject={Second pinching problem for minimal hypersurfaces},pdfauthor={Jinxi Zhang}}
\begin{document}
\begin{abstract}
Let $M^4$ be a closed connected manifold minimally immersed in
$\Sph^5(1)$, and let $h$ be its second fundamental form. We prove
that constant $S=|h|^2>4$ implies $S>8$. No constancy of
$\tr(h^3)$ is assumed. The proof combines integral curvature identities
with a joint estimate for two contractions of the symmetrized Hessian.
The remaining algebraic inequality is verified in Mathematica using
exact rational arithmetic. The code and data are included in the
supplementary material.
\end{abstract}
\maketitle

\section{Introduction}
Let $F:M^n\to\Sph^{n+1}(1)$ be a smooth minimal immersion of a closed
connected manifold, with second fundamental form $h$ and $S=|h|^2$.
Here closed means compact without boundary. The Gauss equation
\[
 \Scal_M=n(n-1)-S
\]
shows that the scalar curvature $\Scal_M$ of the induced metric
is constant if and only if $S$ is constant. Let
$\mathcal C_n$ denote the set of constant values of $S$ realized
by such immersions.

\begin{conjecture}[Chern; see \cite{CDK,HXZ}]
For each $n\ge2$, the set $\mathcal C_n$ is discrete.
\end{conjecture}

The \emph{strong Chern conjecture} asserts that every such immersion
with constant $S$ is isoparametric: its principal curvatures and their
multiplicities are locally constant; see \cite[Conjecture~1.4]{HXZ}.
In dimension four, the minimal isoparametric examples have $S=0,4,12$
\cite{Cartan,Munzner1,Munzner2}.
The \emph{second pinching problem} asks whether $S>n$ implies
$S\ge2n$ when $S$ is constant. We prove a strict bound in dimension four.

\begin{theorem}\label{thm:main}
Let $F:M^4\to\Sph^5(1)$ be a smooth minimal immersion of a nonempty
closed connected manifold, and let $h$ be its second fundamental form.
If $S=|h|^2$ is constant and $S>4$, then
\[
 S>8.
\]
\end{theorem}

Recall the classical first pinching theorem.

\begin{theorem}[\cite{Simons,Lawson,CDK}]\label{thm:first}
If $M^n$ is closed and connected, $F:M^n\to\Sph^{n+1}(1)$ is minimal,
and $S=|h|^2\le n$, then either $S\equiv0$ or $S\equiv n$. The image is,
respectively, an equatorial sphere or a Clifford minimal hypersurface
\[
 \Sph^k\!\left(\sqrt{k/n}\right)\times
 \Sph^{n-k}\!\left(\sqrt{(n-k)/n}\right),\qquad 1\le k\le n-1.
\]
\end{theorem}

Together, these theorems give the following classification for $S\le8$.

\begin{corollary}\label{cor:main}
Let $F:M^4\to\Sph^5(1)$ be a smooth minimal immersion of a nonempty
closed connected manifold. Let $h$ be its second fundamental form.
If $S=|h|^2$ is constant and $S\le8$, then $S=0$ or $S=4$.
Up to an ambient isometry, the image
is an equatorial $\Sph^4$, or one of
\[
 \Sph^1(1/2)\times\Sph^3(\sqrt3/2),\qquad
 \Sph^2(1/\sqrt2)\times\Sph^2(1/\sqrt2).
\]
In particular, $\mathcal C_4\cap[0,8]=\{0,4\}$.
\end{corollary}
\begin{proof}
If $S>4$, Theorem~\ref{thm:main} contradicts $S\le8$. Apply
Theorem~\ref{thm:first} to the remaining interval.
\end{proof}
The region $S>8$ remains unclassified here.

For earlier bounds on the second gap, see \cite{PT1,PT2,SY}. For closed connected
minimal immersions with constant $S>n$, Chen \cite[Theorem~1.1]{Chen}
obtains $S>85n/59$ when $n\ge4$. For complete minimal immersions in
dimension four with constant $S>4$, Cheng--Li--Wei
\cite[Theorem~1.2]{CLW} obtain $S>200000/29609$.
Neither result assumes constant $f_3=\tr(h^3)$.

For closed connected minimal four-dimensional immersions with constant
$S$, isoparametricity follows from constant third mean curvature
$H_3=f_3/12$ \cite[Theorem~1.8]{HXZ} or constant Gauss--Kronecker
curvature \cite[Theorem~1.1]{GLLY}. With constant $S,f_3$, the implication
$S>n\Rightarrow S\ge2n$ holds for closed connected minimal immersions
with $n\ge5$ \cite[Theorem~1.1]{GTYZ}, and for closed oriented minimal
immersions with $n\ge3$ \cite[Theorem~1.1]{Guan}.

\subsection*{Proof strategy}
The argument combines pointwise estimates with integral curvature
identities, as in
\cite[Sections~2 and~4]{LXX} and \cite[Section~4]{Chen}.
We use the curvature defect, Hessian symmetrization and two trace
constraints in the form given in \cite[Sections~2.2, 3.1 and~3.3]{Chen}.

The two contraction constraints act on the same reduced Hessian.
A single square-completion estimate retains the mixed term between
the two contractions. The second constraint involves $\nabla f_3$,
represented after normalization by the vector $d$ in \eqref{eq:d}.

Assume $4<S\le8$ and put $t=(S-4)/S$. A joint estimate for two
contractions of the symmetrized Hessian, combined with three
zero-average identities, gives a quadratic polynomial $W_t$ with
nonnegative average. For each fixed admissible pair of normalized
tensors (Definition~\ref{def:admissible}), $W_0\le0$ and $W_t$ is convex in $t$.
Appendix~\ref{app:certificate} proves $W_{1/2}\le-10^{-4}$
by a finite calculation using exact rational arithmetic. Convexity gives
$W_t<0$ throughout $0<t\le1/2$, contradicting its nonnegative average.

\section{Curvature identities and normalization}\label{sec:curvature}
We first collect the curvature identities that determine the Hessian
norm, then normalize the second fundamental form and its first
derivative. The normalization separates the constant $S$ from the
pointwise algebraic constraints.

Let $F:M^4\to\Sph^5(1)$ be a smooth minimal immersion. The metric
and covariant derivative $\nabla$ are induced by $F$. Choose a local
unit normal and an orthonormal tangent frame $e_1,\ldots,e_4$.
In this frame, let $h=(h_{ij})$ be the second fundamental form,
$S=|h|^2$, and $f_j=\tr(h^j)$. Powers of $h$ refer to the
associated self-adjoint endomorphism.

Indices range from $1$ to $4$. For covariant tensors of order $r$,
\[
 \langle U,V\rangle=\sum_{i_1,\ldots,i_r=1}^4
 U_{i_1\cdots i_r}V_{i_1\cdots i_r},\qquad |U|^2=\langle U,U\rangle.
\]
Thus $\langle X,Y\rangle=\tr(XY)$ for symmetric matrices.
Write $I$ for the $4\times4$ identity matrix and $\Sym^r\R^4$ for
the symmetric covariant $r$-tensors. For symmetric matrices,
$X\preceq Y$ means $v^T(Y-X)v\ge0$ for every $v\in\R^4$;
$X\prec Y$ means strict inequality for every $v\ne0$.
The derivative and Laplacian conventions are
\[
 h_{ijk}=(\nabla_k h)_{ij},\qquad
 h_{ijkl}=(\nabla_l\nabla h)_{ijk},\qquad
 \Delta=\diver\nabla.
\]
For a scalar function $f$, write $f_{;k}=(\nabla f)_k$.
On tensors, $\Delta$ is the rough Laplacian with the same sign.
Pointwise formulas use a local unit normal. If $M$ is closed and has
no global unit normal, the integral formulas are applied on its compact
unit normal double cover.

The scalar integrands are unchanged by normal
reversal, so their normalized integrals agree with those on $M$.

Let $\delta_{ij}$ be the Kronecker symbol and $\mathrm{Rm}_{ijkl}$
the curvature tensor. The Gauss, Codazzi and Ricci equations are
\begin{align}
 \mathrm{Rm}_{ijkl}&=\delta_{ik}\delta_{jl}-\delta_{il}\delta_{jk}
              +h_{ik}h_{jl}-h_{il}h_{jk},\label{eq:gauss}\\
 h_{ijk}&=h_{ikj},\label{eq:codazzi}\\
 h_{ijkl}-h_{ijlk}
 &=\sum_r h_{rj}\mathrm{Rm}_{rikl}+\sum_r h_{ir}\mathrm{Rm}_{rjkl}.\label{eq:ricci}
\end{align}
Thus $\nabla h$ is fully symmetric and $\nabla^2h$ is symmetric in its
first three indices. At a fixed point choose a principal frame
$h_{ij}=\lambda_i\delta_{ij}$, where $\lambda_i$ are the principal
curvatures. Covariant differentiations are performed
before this choice. In particular,
\begin{equation}\label{eq:riccipair}
 h_{iijj}-h_{jjii}=(\lambda_i-\lambda_j)(1+\lambda_i\lambda_j).
\end{equation}

Define the three scalar contractions
\begin{equation}\label{eq:unnormalized}
 \mathcal A=\sum_{ijk}\lambda_i^2h_{ijk}^2,\qquad
 \mathcal B=\sum_{ijk}\lambda_i\lambda_jh_{ijk}^2,\qquad
 \mathcal C=\sum_{ijk}\lambda_i h_{ijk}^2.
\end{equation}
\begin{lemma}[Curvature identities]\label{lem:classical}
Let $F:M^4\to\Sph^5(1)$ be a smooth minimal immersion.
Let $h$ be its second fundamental form, $S=|h|^2$,
$f_j=\tr(h^j)$, and let $\mathcal A,\mathcal B,\mathcal C$ be
as in \eqref{eq:unnormalized}. Then
\begin{gather}
 \Delta h=(4-S)h,\qquad
 \tfrac12\Delta S=|\nabla h|^2+S(4-S),\label{eq:simons}\\
 \tfrac12\Delta|\nabla h|^2
 =|\nabla^2h|^2+(11-S)|\nabla h|^2
       -3(\mathcal A-2\mathcal B)-\tfrac32|\nabla S|^2.\label{eq:PT}
\end{gather}
Moreover,
\begin{align}
 \Delta f_3&=3(4-S)f_3+6\mathcal C,
 & (f_3)_{;k}&=3\sum_i\lambda_i^2h_{iik},\label{eq:f3}\\
 \Delta f_4&=4(4-S)f_4+8\mathcal A+4\mathcal B.\label{eq:f4}
\end{align}
\end{lemma}
\begin{proof}
The tensor equation and the first scalar equation in \eqref{eq:simons}
are Simons' identities \cite{Simons}. Equation~\eqref{eq:PT} is the
Peng--Terng identity \cite{PT1,PT2}; the scalar formulas, with these
signs, are recorded in \cite[Section~2, equations~(2.1)--(2.2)]{LXX}.
The product rule applied to $\tr(h^3)$ and $\tr(h^4)$ gives
\eqref{eq:f3} and \eqref{eq:f4}.
\end{proof}

\subsection{Normalized tensors}
From this point on, the geometric applications have constant $S>4$.
The following normalization gives unit tensors and puts the interval
$4<S\le8$ in correspondence with $0<t\le1/2$. Define
\begin{equation}\label{eq:normalization}
 t=\frac{S-4}{S},\qquad A=\frac h{\sqrt S},\qquad
 \mu_i=\frac{\lambda_i}{\sqrt S},\qquad
 \tau=\frac{\nabla h}{\sqrt{S(S-4)}}.
\end{equation}
Simons' identity gives $|\nabla h|^2=tS^2$. Thus $A$ and $\tau$
have norm one.

\begin{definition}[Admissible pair]\label{def:admissible}
An \emph{admissible pair} consists of a symmetric matrix
$A\in\R^{4\times4}$ and a tensor $\tau\in\Sym^3\R^4$ such that
\begin{equation}\label{eq:admissible}
 \tr A=0,\qquad |A|=|\tau|=1,\qquad
 \sum_i\tau_{iik}=0,\qquad
 \sum_{ij}A_{ij}\tau_{ijk}=0\quad(1\le k\le4).
\end{equation}
In a basis with $A=\diag(\mu_1,\ldots,\mu_4)$, the last condition
is $\sum_i\mu_i\tau_{iik}=0$. All component formulas below use
such a basis unless a different basis is specified.
\end{definition}
The pair \eqref{eq:normalization} is admissible: differentiating
$\tr h=0$ gives its ordinary trace condition, and differentiating
$S$ gives its $A$-weighted trace condition.

For each $k$, the symmetric matrix $(\tau_{ijk})_{ij}$ is orthogonal
to $I$ and $A$ by the two trace constraints. Its contraction with
$A^2$ therefore depends only on the component of $A^2$ orthogonal
to $I$ and $A$.
For an admissible pair $(A,\tau)$, define the scalars $m,u,p$ and
the symmetric matrix $\Phi$ by
\begin{align}
 m&=\tr A^3,\qquad u=\tr A^4,\label{eq:moments}\\
 \Phi&=A^2-mA-\tfrac14I,\qquad
 p=|\Phi|^2=u-m^2-\tfrac14.\label{eq:phi}
\end{align}
The scalars $a,b,c$ below are the normalized versions of
$\mathcal A,\mathcal B,\mathcal C$. We also define the contraction
$d$ of $A^2$ with $\tau$ and the Gram matrix $R$ of the four
matrices $(\tau_{ijk})_{ij}$, indexed by $k$:
\begin{align}
 a&=\sum_{ijk}\mu_i^2\tau_{ijk}^2,\qquad
 b=\sum_{ijk}\mu_i\mu_j\tau_{ijk}^2,\qquad
 c=\sum_{ijk}\mu_i\tau_{ijk}^2,\label{eq:abc}\\
 d_k&=\sum_i\mu_i^2\tau_{iik},\label{eq:d}\\
 R_{kl}&=\sum_{ij}\tau_{ijk}\tau_{ijl}.\label{eq:Rmu}
\end{align}
Thus $d\in\R^4$ and $R\in\Sym^2\R^4$, with
$R\succeq0$ and $\tr R=1$.
For the geometric pair \eqref{eq:normalization}, these contractions
recover the unnormalized quantities:
\[
 \mathcal A=tS^3a,\qquad \mathcal B=tS^3b,\qquad
 \mathcal C=tS^{5/2}c.
\]
The vector $d$ records the variation of the cubic trace:
\[
 \nabla f_3=3S^2\sqrt t\,d.
\]
The two trace constraints in \eqref{eq:admissible} do not imply $d=0$.
The identities $\tr A=0$ and $|A|=1$ imply
\begin{equation}\label{eq:phiorth}
 \tr\Phi=\tr(A\Phi)=0,\qquad
 p\ge0,\qquad \mu_i^2\le\tfrac34.
\end{equation}
The last inequality follows from
$\mu_i^2=(\sum_{j\ne i}\mu_j)^2\le3(1-\mu_i^2)$.
Since $S$ and $|\nabla h|^2$ are constant, \eqref{eq:PT} now reads
\begin{equation}\label{eq:PTnorm}
 \frac{|\nabla^2h|^2}{S^3}
       =3t(a-2b)+\frac{t(11t-7)}4.
\end{equation}

\section{The reduced Hessian}\label{sec:Hessian}
Equation~\eqref{eq:PTnorm} gives the normalized squared norm of the
Hessian. Differentiating the constant quantities $\tr h$, $S$ and
$|\nabla h|^2$ gives three of its contractions.
To use this information in a norm estimate, we pass to a fully symmetric
tensor and remove its ordinary trace. The tensor that remains will have
two prescribed contractions and a known squared norm.

\subsection{The reduction statement}
The three contractions are the ordinary trace, the $A$-weighted trace,
and contraction against the cubic tensor $\tau$.

\begin{definition}[Contraction maps]\label{def:maps}
Fix an admissible pair $(A,\tau)$. Define
\[
 \T_0,\T_A:\Sym^4\R^4\longrightarrow\Sym^2\R^4,
 \qquad \N_\tau:\Sym^4\R^4\longrightarrow\R^4
\]
by
\begin{equation}\label{eq:maps}
 (\T_0U)_{kl}=\sum_iU_{iikl},\quad
 (\T_AU)_{kl}=\sum_i\mu_iU_{iikl},\quad
 (\N_\tau U)_l=\sum_{ijk}\tau_{ijk}U_{ijkl}.
\end{equation}
Let
\[
 \Hfour=\{U\in\Sym^4\R^4:\T_0U=0\},
 \qquad \Pz:\Sym^4\R^4\longrightarrow\Hfour
\]
be, respectively, the trace-free subspace and its orthogonal projection.
For four-tensors symmetric only in the first three indices,
$\T_0,\T_A,\N_\tau$ still denote the component contractions in
\eqref{eq:maps}.
\end{definition}

The two contractions retained after trace removal will be denoted by
$E_t$ and $b_0$. Write
\[
 \Szero=\{X\in\R^{4\times4}:X^T=X,\ \tr X=0\}.
\]
For an admissible pair $(A,\tau)$ and $t\in\R$, define
\begin{equation}\label{eq:Edefinitions}
 \Psi=\tfrac14A^2+\tfrac3{16}I-R,\qquad
 E_t=\tfrac12\Phi+t\Psi,\qquad b_0=-\tfrac32Ad.
\end{equation}
Here $\Psi,E_t\in\Szero$ and $b_0\in\R^4$. The term $\Phi/2$ is
independent of $t$, while $\Psi$ is the coefficient of $t$ in $E_t$.

\begin{proposition}[Hessian contractions]\label{prop:reduction}
Let $F:M^4\to\Sph^5(1)$ be a smooth minimal immersion with
constant $S=|h|^2>4$. At each point, define $t,A,\tau$ by
\eqref{eq:normalization} and $E_t,b_0$ by \eqref{eq:Edefinitions}.
There exists $V\in\Hfour$ such that
\begin{align}
 \T_AV&=E_t,\label{eq:Et}\\
 \N_\tau V&=b_0,\label{eq:b0}\\
 |V|^2&=3t(a-2b)+\frac{t(35t-28)}{16}-\tfrac32p.\label{eq:Vnorm}
\end{align}
\end{proposition}

\subsection{Symmetrization and trace removal}
We prove the proposition by subtracting two orthogonal corrections:
the part fixed by the Ricci equation, and the least-norm tensor with
the remaining ordinary trace. The latter correction is supplied by
Lemma~\ref{lem:T0}.

For a covariant $r$-tensor $T$, define
\[
 (\Sym T)_{i_1\cdots i_r}
 =\frac1{r!}\sum_{\sigma\in\mathfrak S_r}
 T_{i_{\sigma(1)}\cdots i_{\sigma(r)}},
\]
where $\mathfrak S_r$ is the permutation group.
Adjoints $L^*$ are taken with respect to the ordered-index inner product
of Section~\ref{sec:curvature}. The product $F\otimes G$ denotes the
tensor product without symmetrization.

\begin{lemma}[Ordinary trace and its adjoint]\label{lem:T0}
For every $X\in\Sym^2\R^4$,
\begin{gather}
 \T_0^*X=\Sym(I\otimes X),\qquad
 \T_0\T_0^*X=\tfrac43X+\tfrac16\tr(X)I,\label{eq:T0Gram}\\
 (\T_0\T_0^*)^{-1}X=\tfrac34X-\tfrac1{16}\tr(X)I,\qquad
 \Pz=\operatorname{Id}_{\Sym^4\R^4}-\T_0^*(\T_0\T_0^*)^{-1}\T_0.\label{eq:P0}
\end{gather}
\end{lemma}
\begin{proof}
The adjoint formula follows by symmetry. In components,
\begin{align*}
 (\T_0^*X)_{ijkl}=\tfrac16(&\delta_{ij}X_{kl}+\delta_{ik}X_{jl}
 +\delta_{il}X_{jk}\\
 &+\delta_{jk}X_{il}+\delta_{jl}X_{ik}+\delta_{kl}X_{ij}).
\end{align*}
Tracing the first two indices proves \eqref{eq:T0Gram}.
This operator acts by $4/3$ on trace-free matrices and by $2$ on
multiples of $I$, giving the stated inverse.
In \eqref{eq:P0}, the subtracted tensor has the same ordinary trace
as the input and lies in $\operatorname{im}\T_0^*=(\ker\T_0)^\perp$.
The remainder is therefore the orthogonal projection onto $\Hfour$.
\end{proof}

\begin{proof}[Proof of Proposition~\ref{prop:reduction}]
Set $\mathcal D=S^{-3/2}\nabla^2h$. Since $S$ is constant,
\eqref{eq:simons} makes $|\nabla h|^2$ constant as well.
Differentiating $\tr h=0$ and $S$ twice, and $|\nabla h|^2$ once, gives
\begin{equation}\label{eq:rawtraces}
 \T_0\mathcal D=0,\qquad
 \T_A\mathcal D=-tR,\qquad
 \N_\tau\mathcal D=0.
\end{equation}
First remove the failure of full symmetry. Put
$\widehat U=\Sym\mathcal D$, $C=\mathcal D-\widehat U$ and
$\kappa=S^{-1}=(1-t)/4$.
Since $\mathcal D$ is symmetric in its first three indices,
\[
 \widehat U_{ijkl}=\tfrac14(\mathcal D_{ijkl}+\mathcal D_{ijlk}
                         +\mathcal D_{iklj}+\mathcal D_{jkli}).
\]
\Needspace{4\baselineskip}
For $i\ne j$, symmetry in the first three indices gives
$\widehat U_{iijj}=(\mathcal D_{iijj}+\mathcal D_{jjii})/2$.
Thus $C_{iijj}=(\mathcal D_{iijj}-\mathcal D_{jjii})/2$;
\eqref{eq:riccipair}, divided by $S^{3/2}$, determines its sign
and the factor $1/2$. In a principal frame, \eqref{eq:gauss}
and~\eqref{eq:ricci} show that $C$ vanishes outside the index patterns
$\{i,i,j,j\}$.
The possibly nonzero components are therefore
\begin{gather*}
 C_{iijj}=\tfrac12(\mu_i-\mu_j)(\kappa+\mu_i\mu_j),\qquad
 C_{ijij}=C_{jiij}=C_{iijj},\\*
 C_{jjii}=C_{jiji}=C_{ijji}=-C_{iijj}.
\end{gather*}

Summing these components, using $\sum_i\mu_i=0$, $\sum_i\mu_i^2=1$
and $\sum_i\mu_i^3=m$, gives the first two contractions below.
For the third, each pattern $(i,i,l)$ with $i\ne l$ occurs three
times in the ordered contraction, giving
\begin{align}
 \T_0C&=\tfrac t2A,&
 \T_AC&=\tfrac12(\kappa I+mA-A^2),\label{eq:Ctraces}\\
 (\N_\tau C)_l
 &=\tfrac32\sum_i(\mu_i-\mu_l)(\kappa+\mu_i\mu_l)\tau_{iil}
 =\tfrac32\mu_ld_l.\label{eq:Cthird}
\end{align}
Indeed, the weight in this sum is
$(\kappa-\mu_l^2)\mu_i-\kappa\mu_l+\mu_l\mu_i^2$.
The first two terms vanish after summation by the two trace constraints;
the last gives $\mu_ld_l$ by \eqref{eq:d}.
Subtracting these contractions from \eqref{eq:rawtraces} gives
\begin{equation}\label{eq:Utraces}
 \T_0\widehat U=-\tfrac t2A,\qquad
 \T_A\widehat U=-tR-\tfrac12(\kappa I+mA-A^2),\qquad
 \N_\tau\widehat U=-\tfrac32Ad.
\end{equation}

Next remove the ordinary trace of $\widehat U$.
By Lemma~\ref{lem:T0}, the least-norm tensor with trace
$-tA/2$ is
\[
 U_0=\T_0^*(\T_0\T_0^*)^{-1}(-tA/2)
     =-\tfrac{3t}{8}\Sym(I\otimes A).
\]
Its squared norm follows from
$|U_0|^2=\langle-tA/2,(\T_0\T_0^*)^{-1}(-tA/2)\rangle$;
its contractions follow from the six-term adjoint formula in
Lemma~\ref{lem:T0}. Thus
\begin{equation}\label{eq:U0}
 |U_0|^2=\tfrac3{16}t^2,\qquad
 \T_AU_0=-\tfrac t{16}(I+4A^2),\qquad \N_\tau U_0=0.
\end{equation}
Set $V=\widehat U-U_0$. Then $V\in\Hfour$ and
\eqref{eq:Utraces}--\eqref{eq:U0} give \eqref{eq:Et} and
\eqref{eq:b0}.

It remains to compute $|V|^2$. Symmetrization is an orthogonal projection, so
$C\perp\Sym^4\R^4$; also $U_0\in(\ker\T_0)^\perp$ and
$V\in\ker\T_0$. Thus $\mathcal D=C+U_0+V$ is an orthogonal
decomposition.

The components of $C$ give
\begin{align}
 |C|^2
 &=\tfrac34\sum_{ij}(\mu_i-\mu_j)^2(\kappa+\mu_i\mu_j)^2
 =\tfrac32p+\tfrac38t^2.\label{eq:Cnorm}
\end{align}
For each unordered pair $\{i,j\}$, the six ordered components displayed
above have squared sum $6C_{iijj}^2$; this gives the factor $3/4$ in
the sum over ordered pairs. Expanding the sum before multiplication
by $3/4$, and using
$\sum_i\mu_i=0$, $\sum_i\mu_i^2=1$ gives
$8\kappa^2-4\kappa+2u-2m^2=2p+t^2/2$.
Together with $|U_0|^2$ in \eqref{eq:U0}, this gives
\begin{equation}\label{eq:orthnorm}
 \frac{|\nabla^2h|^2}{S^3}=|V|^2+\tfrac32p+\tfrac9{16}t^2.
\end{equation}
Equation~\eqref{eq:PTnorm} now gives
\eqref{eq:Vnorm}.
\end{proof}

\section{The joint contraction estimate}\label{sec:algebra}
We now bound the squared norm of a tensor in $\Hfour$ from its two
prescribed contractions. The estimate is purely algebraic. For the
geometric application in Proposition~\ref{prop:reduction}, take
$E=E_t$ and $b_0=-\tfrac32Ad$.
Fix an admissible pair $(A,\tau)$.

\subsection{The estimate}
Suppose $V\in\Hfour$ satisfies $\T_AV=E$ and $\N_\tau V=b_0$,
where $E\in\Szero$ and $b_0\in\R^4$. For any $Y\in\Szero$ and
$v\in\R^4$, square completion gives
\[
 |V|^2\ge2\langle E,Y\rangle+2b_0\cdot v
          -|\Pz(\T_A^*Y+\N_\tau^*v)|^2,
\]
where $b_0\cdot v=b_0^Tv$. The two terms inside the projection
contribute a cross term. The squared norm has a matrix term, a
mixed term and a vector term, represented by $\K_A$, $Z$ and
$H_\tau$, respectively.

For $Y\in\Szero$, set
\begin{align}
 \K_AY={}&\tfrac16Y+\tfrac14(A^2Y+YA^2)-\tfrac16AYA
 \notag\\
 &+\tfrac7{36}\tr(AY)A-\tfrac1{12}\tr(A^2Y)I,\label{eq:K}\\
 (ZY)_l={}&\tfrac18\sum_{ij}(4\mu_l-\mu_i-\mu_j)\tau_{ijl}Y_{ij},
 \qquad H_\tau=\tfrac14I+\tfrac9{16}R.\label{eq:ZH}
\end{align}
The map $Z$ depends on both $A$ and $\tau$. The domains and codomains are
\[
 \K_A:\Szero\to\Szero,\qquad
 Z:\Szero\to\R^4,\qquad H_\tau:\R^4\to\R^4.
\]

\begin{proposition}[Two prescribed contractions]\label{prop:joint}
Let $(A,\tau)$ be admissible. Let $E\in\Szero$, $b_0\in\R^4$,
and $V\in\Hfour$ satisfy
\[
 \T_AV=E,\qquad \N_\tau V=b_0.
\]
For every $Y\in\Szero$ and $v\in\R^4$,
\begin{equation}\label{eq:joint}
 |V|^2\ge2\langle E,Y\rangle+2b_0\cdot v
 -\langle Y,\K_AY\rangle-2v\cdot ZY-v^TH_\tau v.
\end{equation}
\end{proposition}

\subsection{The Gram identities}
To prove Proposition~\ref{prop:joint}, we identify the three terms
in the projected squared norm. The same calculation gives a lower
bound for $\K_A$ that will control the quadratic coefficient in
Section~\ref{sec:integral}.

\begin{lemma}[Contraction identities]\label{lem:Gram}
For an admissible pair $(A,\tau)$,
\begin{equation}\label{eq:Grams}
 \T_A\Pz\T_A^*=\K_A,\qquad
 \N_\tau\Pz\T_A^*=Z,\qquad
 \N_\tau\Pz\N_\tau^*=H_\tau.
\end{equation}
The first two identities act on $\Szero$, and the third on $\R^4$.
For every $Y\in\Szero$,
\begin{equation}\label{eq:Klower}
 \langle Y,\K_AY\rangle
 \ge\tfrac16|Y|^2+\tfrac13\tr(A^2Y^2).
\end{equation}
\end{lemma}
\begin{proof}
The adjoints are $\T_A^*Y=\Sym(A\otimes Y)$ and
$\N_\tau^*v=\Sym(\tau\otimes v)$.
For the matrix term, compute the unprojected contraction and then
subtract the ordinary trace correction from \eqref{eq:P0}:
\begin{align*}
 \T_A\T_A^*Y
 &=\tfrac16Y+\tfrac13(A^2Y+YA^2)+\tfrac16\tr(AY)A,\\
 \T_0\T_A^*Y&=\tfrac13(AY+YA),\\
 \T_A\T_0^*(\T_0\T_0^*)^{-1}\T_0\T_A^*Y
 &=\tfrac1{12}(A^2Y+2AYA+YA^2)-\tfrac1{36}\tr(AY)A
       +\tfrac1{12}\tr(A^2Y)I.
\end{align*}
The first expression minus the third is $\K_AY$, proving the
matrix identity.

For the mixed term, the vanishing $A$-weighted trace of $\tau$ gives
\[
 (\N_\tau\T_A^*Y)_l=\tfrac12\mu_l\sum_{ij}\tau_{ijl}Y_{ij}.
\]
Substitution in \eqref{eq:P0} gives the trace correction
$\frac18\sum_{ij}(\mu_i+\mu_j)\tau_{ijl}Y_{ij}$; the scalar part
of $(\T_0\T_0^*)^{-1}$ contributes nothing because the ordinary trace
of $\tau$ vanishes. Subtraction gives $Z$.

For the vector term,
\[
 \N_\tau\N_\tau^*=\tfrac14I+\tfrac34R,\qquad
 (\T_0\N_\tau^*v)_{ij}=\tfrac12\sum_k\tau_{ijk}v_k.
\]
The latter matrix is trace-free and has squared norm $v^TRv/4$.
On trace-free matrices, $(\T_0\T_0^*)^{-1}$ acts by $3/4$.
The projection therefore subtracts $3R/16$, giving $H_\tau$.

To obtain \eqref{eq:Klower}, expand
\[
 \langle Y,\K_AY\rangle
 =\tfrac16|Y|^2+\tfrac12\tr(A^2Y^2)
   -\tfrac16\tr(AYAY)+\tfrac7{36}\tr(AY)^2
\]
The squared trace term is nonnegative, and
\[
 \tr(A^2Y^2)-\tr(AYAY)
 =\frac12\sum_{ij}(\mu_i-\mu_j)^2Y_{ij}^2\ge0.\qedhere
\]
\end{proof}

\begin{proof}[Proof of Proposition~\ref{prop:joint}]
Set $U=\Pz(\T_A^*Y+\N_\tau^*v)$. Since $V=\Pz V$,
\[
 \langle V,U\rangle=\langle E,Y\rangle+b_0\cdot v.
\]
Lemma~\ref{lem:Gram} gives
\[
 |U|^2=\langle Y,\K_AY\rangle+2v\cdot ZY+v^TH_\tau v.
\]
Substitution in $0\le|V-U|^2=|V|^2-2\langle V,U\rangle+|U|^2$
proves \eqref{eq:joint}.
\end{proof}

\section{The averaged inequality and the contradiction}\label{sec:integral}
For the geometric tensor in Proposition~\ref{prop:reduction}, the
difference between its squared norm and the lower bound in
Proposition~\ref{prop:joint} is itself a squared norm. We will add
three divergence terms to this difference, preserving its average.
The resulting polynomial $W_t$ will be quadratic in $t$; convexity
will reduce its sign on $0<t\le1/2$ to the two endpoints.

For a nonempty closed Riemannian manifold $M$ and an integrable
scalar function $g$, write
\[
 \av g=\frac1{\Vol(M)}\int_M g\,d\Vol.
\]

\subsection{The polynomial \texorpdfstring{$W_t$}{W(t)}}
We use \eqref{eq:joint} with $E=E_t$. Its terms linear in $Y$ are
$2\langle E_t-Z^*v,Y\rangle$. We take $v$ proportional to $b_0$
and $Y_t$ linear in $t$, with coefficient proportional to the matrix
$E_{1/2}-Z^*v$.

For an admissible pair $(A,\tau)$ and real $t$, define
\begin{equation}\label{eq:L}
 L_4=t(2a+b-u),\qquad
 L_3=t(-m^2+2mc+3|d|^2),\qquad
 J=t(a-2b)-p-\tfrac t4.
\end{equation}
On the geometric pair these are divergences; their zero averages
will be proved in Lemma~\ref{lem:mean}.

\begin{definition}[The polynomial $W_t$]\label{def:Wt}
Fix an admissible pair $(A,\tau)$.
For a real parameter $t$, set
\begin{equation}\label{eq:choice}
 v=-7Ad,\qquad B_t=E_t-Z^*v,\qquad
 B=B_{1/2},\qquad Y_t=8tB.
\end{equation}
Thus $v\in\R^4$, $B_t,B,Y_t\in\Szero$, and $Y_{1/2}=4B$.
The pair $(A,\tau)$ and the matrix $B$ are held fixed as $t$ varies.
Let $Q_t$ be the right-hand side of \eqref{eq:joint} for these choices.
Using the norm expression \eqref{eq:Vnorm} and $L_4,L_3,J$ from
\eqref{eq:L}, set
\begin{align}
 Q_t={}&16t\langle B_t,B\rangle-64t^2\langle B,\K_AB\rangle
               +2b_0\cdot v-v^TH_\tau v,\label{eq:Qt}\\
 W_t={}&3t(a-2b)+\frac{t(35t-28)}{16}-\tfrac32p-Q_t
              -\tfrac{89}{8}L_4+16L_3+\tfrac{13}{4}J.\label{eq:Wt}
\end{align}

The adjoint $Z^*:\R^4\to\Szero$ and $B$ have components
\begin{equation}\label{eq:Zadjoint}
 \begin{aligned}
 (Z^*v)_{ij}&=\tfrac18\sum_l(4\mu_l-\mu_i-\mu_j)\tau_{ijl}v_l,\\
 B&=\tfrac58A^2-\tfrac m2A-\tfrac1{32}I-\tfrac12R+7Z^*Ad.
 \end{aligned}
\end{equation}
\end{definition}

The matrix trial vanishes at $t=0$ and is proportional to $B_{1/2}$
at $t=1/2$. Keeping $B=B_{1/2}$ fixed makes $W_t$ quadratic in $t$;
the divergence terms change only its constant and linear coefficients.
The constants in \eqref{eq:choice} and~\eqref{eq:Wt} must give
nonpositivity at $t=0$, convexity in $t$, and strict negativity at $t=1/2$.
Proposition~\ref{prop:convex} proves the first two; Appendix~\ref{app:certificate}
proves the third using positive-definite Gram representations checked
in exact rational arithmetic.

\subsection{The nonnegative average}
The formulas for $a,b,c,d,R$ are tensor contractions of $A$ and
$\tau$, and \eqref{eq:Grams} expresses $\K_A,Z,H_\tau$ in terms
of contractions and orthogonal projection. Thus $W_t$ is independent
of the orthonormal frame. Reversing the unit normal sends
$(A,\tau)$ to $(-A,-\tau)$ and also leaves $W_t$ unchanged.
For the geometric pair, $W_t$ therefore defines a smooth scalar
function on $M$, including at repeated principal curvatures.

The joint estimate gives a nonnegative average for $W_t$ once
$L_4,L_3,J$ are shown to have zero average.

\begin{proposition}[The integral inequality]\label{prop:meanconvex}
Let $F:M^4\to\Sph^5(1)$ be a smooth minimal immersion of a nonempty
closed manifold, with constant $S=|h|^2>4$. Define $t,A,\tau$ by
\eqref{eq:normalization} and $W_t$ by Definition~\ref{def:Wt}.
Then
\begin{equation}\label{eq:Wmean}
 \av{W_t}\ge0.
\end{equation}
\end{proposition}

The following divergence identities supply these zero averages.

\begin{lemma}[Three zero averages]\label{lem:mean}
Let $F:M^4\to\Sph^5(1)$ be a smooth minimal immersion of a nonempty
closed manifold, with constant $S=|h|^2>4$. Define $t,A,\tau$ by
\eqref{eq:normalization}, and $L_4,L_3,J$ by \eqref{eq:L}. Then
\begin{equation}\label{eq:divergences}
 L_4=\frac{\Delta(Sf_4/4)}{S^4},\qquad
 L_3=\frac{\Delta(f_3^2/6)}{S^4},\qquad
 J=-\frac1{S^3}\diver(h\nabla f_3/3).
\end{equation}
Consequently,
\[
 \av{L_4}=\av{L_3}=\av J=0.
\]
\end{lemma}
\begin{proof}
Equation~\eqref{eq:f4} gives the first formula. For the second, use
\[
 \Delta(f_3^2/6)=\tfrac13f_3\Delta f_3+\tfrac13|\nabla f_3|^2,
 \qquad |\nabla f_3|^2=9tS^4|d|^2.
\]
For the third identity, put $D=\nabla f_3/3$ and compute $\diver(hD)$.
In a principal frame, $D_k=\sum_i\lambda_i^2h_{iik}=(f_3)_{;k}/3$.
Codazzi and $\tr h=0$ give $\diver h=0$, so
\begin{align*}
 \diver(hD)
 &=2\mathcal B+\sum_{ik}\lambda_i^2\lambda_k h_{iikk}\\
 &=2\mathcal B-\mathcal A
   +\sum_{ik}\lambda_i^2\lambda_k(\lambda_i-\lambda_k)
                                  (1+\lambda_i\lambda_k)\\
 &=2\mathcal B-\mathcal A+Sf_4-f_3^2-S^2.
\end{align*}
In the second line we used \eqref{eq:riccipair} and
$\sum_k\lambda_kh_{kkii}=-\sum_{jk}h_{jki}^2$, obtained by
differentiating $S$ twice. Division by $-S^3$ gives the stated $J$.
The integral identities follow from the divergence theorem.
\end{proof}

\begin{proof}[Proof of Proposition~\ref{prop:meanconvex}]
Take $V$ from Proposition~\ref{prop:reduction} and put
$U_t=\Pz(\T_A^*Y_t+\N_\tau^*v)$. The adjoint relations and
Lemma~\ref{lem:Gram} give
\[
 |V|^2-Q_t=|V-U_t|^2.
\]
By \eqref{eq:Vnorm}, $W_t$ differs from $|V|^2-Q_t$ by the
linear combination of $L_4,L_3,J$ in \eqref{eq:Wt}.
This combination has zero average by Lemma~\ref{lem:mean}, so
\[
 \av{W_t}=\av{|V-U_t|^2}\ge0.\qedhere
\]
\end{proof}

\subsection{The value at zero and convexity}
It remains to make $W_t$ pointwise negative on $0<t\le1/2$.
We first control its value at zero and its quadratic coefficient.
The evaluation at $t=0$ is algebraic: the geometric
normalization of $\tau$ requires $S>4$.

\begin{proposition}[Convexity in $t$]\label{prop:convex}
Fix an admissible pair $(A,\tau)$, and let $W_t$ be the polynomial
in Definition~\ref{def:Wt}. With $p=|A^2-mA-I/4|^2$ and
$m=\tr A^3$, one has
\begin{equation}\label{eq:Wbounds}
 W_0\le-\tfrac{75}{32}p\le0,\qquad
 \frac{d^2}{dt^2}W_t\ge\tfrac13\quad(t\in\R).
\end{equation}
\end{proposition}

The value at zero is controlled by $H_\tau$ and $Ad$.
The quadratic coefficient is controlled by $\Psi$ and $\K_A$.
The next two lemmas give these bounds.

\begin{lemma}[Bounds for $R$ and $Ad$]\label{lem:bounds}
For an admissible pair $(A,\tau)$,
\begin{equation}\label{eq:bounds}
 R\preceq\tfrac23I,\qquad H_\tau\preceq\tfrac58I,
 \qquad |Ad|^2\le\tfrac p4.
\end{equation}
\end{lemma}
\begin{proof}
To estimate $R$ in a unit direction, choose that direction as $e_1$
and split $\tau$ into
\[
 a_0=\tau_{111},\quad \xi_\alpha=\tau_{11\alpha},\quad
 \Gamma_{\alpha\beta}=\tau_{1\alpha\beta},\quad
 \Xi_{\alpha\beta\gamma}=\tau_{\alpha\beta\gamma},
 \qquad \alpha,\beta,\gamma>1.
\]
Then $\tr\Gamma=-a_0$ and
\[
 |\tau|^2=a_0^2+3|\xi|^2+3|\Gamma|^2+|\Xi|^2,\qquad
 R_{11}=a_0^2+2|\xi|^2+|\Gamma|^2.
\]
Cauchy--Schwarz gives $a_0^2=(\tr\Gamma)^2\le3|\Gamma|^2$.
Consequently
\[
 2|\tau|^2-3R_{11}=-a_0^2+3|\Gamma|^2+2|\Xi|^2\ge0.
\]
This holds in every direction and proves the first two bounds.

For the bound on $Ad$, choose a basis with $A$ diagonal. Fix $k$ and set
$z^{(k)}=(\tau_{iik})_i$ and
$\theta_k=3|z^{(k)}|^2-2\tau_{kkk}^2$.
Here $\theta_k$ is the contribution of the index patterns
$\{i,i,k\}$ to the ordered norm: each pattern with $i\ne k$ occurs
three times, while $(k,k,k)$ occurs once.
Write $\boldsymbol1=(1,1,1,1)^T$ and
$\mu=(\mu_1,\ldots,\mu_4)^T$. The two trace conditions give
$z^{(k)}\perp\boldsymbol1,\mu$.

The projection of $e_k$ onto this common orthogonal complement has
squared norm $3/4-\mu_k^2$. Hence
$\tau_{kkk}^2\le(3/4-\mu_k^2)|z^{(k)}|^2$.
Using $d_k=\sum_i\Phi_{ii}z_i^{(k)}$ and Cauchy--Schwarz gives
\[
 \theta_k\ge(\tfrac32+2\mu_k^2)|z^{(k)}|^2,
 \qquad
 d_k^2=\left(\sum_i\Phi_{ii}z_i^{(k)}\right)^2
       \le\frac{2p\theta_k}{3+4\mu_k^2}.
\]
The repeated-index groups have disjoint ordered supports, so
$\sum_k\theta_k\le1$. Since $\mu_k^2\le3/4$,
\[
 |Ad|^2\le p\sum_k\frac{2\mu_k^2\theta_k}{3+4\mu_k^2}
       \le\tfrac p4.\qedhere
\]

\end{proof}

The next estimate will be applied with $X=8B$ in the quadratic coefficient.

\begin{lemma}[A quadratic bound]\label{lem:Ebound}
For an admissible pair $(A,\tau)$ and every $X\in\Szero$,
\begin{equation}\label{eq:Ebound}
 2\langle \Psi,X\rangle-\langle X,\K_AX\rangle\le\tfrac{97}{48}.
\end{equation}
\end{lemma}
\begin{proof}
By Lemma~\ref{lem:bounds}, the eigenvalues of $R$ lie in $[0,2/3]$
and sum to one. The vertices of this constraint set are the permutations
of $(2/3,1/3,0,0)$, so convexity of the sum of squares gives
$\tr R^2\le5/9$ and $|R-I/4|^2\le11/36$.

The contribution involving $A^2$ is controlled by
\[
 X/2-X^2/3=3I/16-\tfrac13(X-3I/4)^2\preceq3I/16.
\]
Taking the trace against the positive semidefinite matrix $A^2$
preserves this inequality; $\tr A^2=1$ bounds that contribution by
$3/16$. Combining this with \eqref{eq:Klower} gives
\begin{align*}
 2\langle \Psi,X\rangle-\langle X,\K_AX\rangle
 &\le-2\langle R-I/4,X\rangle-\tfrac16|X|^2
       +\tr\bigl(A^2(X/2-X^2/3)\bigr)\\
 &\le-2\langle R-I/4,X\rangle-\tfrac16|X|^2+\tfrac3{16}\\
 &\le6|R-I/4|^2+\tfrac3{16}\le\tfrac{97}{48}.
\end{align*}
The last line completes the square in $X$ and uses the bound on
$|R-I/4|^2$.
\end{proof}

\begin{proof}[Proof of Proposition~\ref{prop:convex}]
For the constant term, put $w=Ad$ and use Lemma~\ref{lem:bounds}:
\[
 W_0=-\tfrac{19}{4}p-21|w|^2+49w^TH_\tau w
     \le-\tfrac{19}{4}p+\tfrac{77}{8}|w|^2
     \le-\tfrac{75}{32}p.
\]
\begin{samepage}
For the quadratic term, $B$ is fixed, $B_t$ is affine in $t$,
and so are $L_4,L_3,J$. Equations~\eqref{eq:Qt}--\eqref{eq:Wt}
therefore give the coefficient of $t^2$ as
\[
 \tfrac{35}{16}+64\langle B,\K_AB\rangle
                    -16\langle \Psi,B\rangle.
\]
Apply \eqref{eq:Ebound} with $X=8B$ to bound this below by
$35/16-97/48=1/6$. Thus $W_t''\ge1/3$.
\end{samepage}
\end{proof}

\subsection{Completing the proof}
Appendix~\ref{app:certificate} proves the following bound for every
admissible pair.

\begin{proposition}[The bound at $t=1/2$]\label{prop:endpoint}
Let $(A,\tau)$ be an admissible pair as in
Definition~\ref{def:admissible}, and let $W_t$ be defined by
Definition~\ref{def:Wt}. Then
\begin{equation}\label{eq:endpointbound}
 W_{1/2}(A,\tau)\le-\frac1{10000}.
\end{equation}
\end{proposition}

We finish by applying convexity on $[0,1/2]$.

\begin{proof}[Proof of Theorem~\ref{thm:main}]
Suppose $4<S\le8$. Then $0<t\le1/2$, and the tensors defined by
\eqref{eq:normalization} satisfy \eqref{eq:admissible} at every point.
Propositions~\ref{prop:convex} and~\ref{prop:endpoint} give
\[
 W_t\le(1-2t)W_0+2tW_{1/2}\le-\frac{t}{5000}<0.
\]
Since $t$ is constant, integration contradicts \eqref{eq:Wmean}.
\end{proof}

\clearpage
\appendix
\section{The bound at \texorpdfstring{$t=1/2$}{t=1/2}}\label{app:certificate}
We prove Proposition~\ref{prop:endpoint} by parametrizing the spectrum
of $A$ and the components of $\tau$, then reducing the inequality to
positivity of a quartic. Bernstein expansion on a rectangular cover of
$[0,1]^2$ reduces this to Gram identities and integer matrix inequalities.
Setting $t=1/2$ in \eqref{eq:Wt} gives
\begin{equation}\label{eq:endpointexpanded}
\begin{split}
 W_{1/2}={}&\tfrac{21}{32}+\tfrac{13}{16}p-\tfrac{39}{16}m^2
 -8a-\tfrac{189}{16}b+16mc+24|d|^2\\
 &-8|B|^2+16\langle B,\K_AB\rangle
 -21|Ad|^2+49(Ad)^TH_\tau(Ad).
\end{split}
\end{equation}

\subsection*{Parametrization}
Since $W_{1/2}$ is invariant under orthogonal changes of basis, we may
diagonalize $A$. Its eigenvalue vector
$\mu=(\mu_1,\ldots,\mu_4)^T$ is a unit vector in
$\boldsymbol1^\perp$, where $\boldsymbol1=(1,1,1,1)^T$.
The separate transformations $A\mapsto-A$ and $\tau\mapsto-\tau$
preserve \eqref{eq:endpointexpanded}. Define
\[
 U=\frac12\begin{pmatrix}
 1&1&1\\ 1&-1&-1\\ -1&1&-1\\ -1&-1&1
 \end{pmatrix}.
\]
Its columns form an orthonormal basis of $\boldsymbol1^\perp$.
Permuting the four eigenvalues acts on $U^T\mu$ by coordinate
permutations and even sign changes. Together with $\mu\mapsto-\mu$,
these allow us to arrange the three coordinates of $U^T\mu$ in
nonincreasing order, all nonnegative. The first coordinate is positive;
dividing by it gives $(1,r,ry)^T$ with $0\le r,y\le1$. If the second
coordinate of $U^T\mu$ is zero, take $r=y=0$. Thus every spectrum is
represented, up to these symmetries, by
\begin{equation}\label{eq:spectralchart}
 \begin{gathered}
 \ell=(1+r+ry,\ 1-r-ry,\ -1+r-ry,\ -1-r+ry),\\
 L=|\ell|^2=4(1+r^2+r^2y^2),\qquad \mu=\ell/\sqrt L.
 \end{gathered}
\end{equation}
For each repeated-index group, the two trace constraints leave the
two-dimensional space $\{\boldsymbol1,\ell\}^\perp$. A basis valid
throughout the closed square is given by the columns of
\begin{equation}\label{eq:cubicchart}
 C(r,y)=\begin{pmatrix}
 -1+ry&-r+ry\\
 1+r&0\\
 0&1+r\\
 -r-ry&-1-ry
 \end{pmatrix}.
\end{equation}
The columns satisfy both trace constraints; the two rows containing
$1+r$ prove independence, including at the boundary. The four
repeated-index groups and four distinct-index components give twelve
coordinates. For $x=(x_1,\ldots,x_{12})\in\R^{12}$, define the symmetric
cubic tensor $T=T(r,y,x)$ by
\[
 (T_{234},T_{134},T_{124},T_{123})=(x_1,x_2,x_3,x_4)
\]
and, for $k=1,\ldots,4$,
\begin{equation}\label{eq:repeatedchart}
 (T_{iik})_{i=1}^4=C(r,y)(x_{2k+3},x_{2k+4})^T.
\end{equation}
Since these groups of components are disjoint, the formulas parametrize
the whole constrained cubic space, with
\[
 N(r,y,x):=|T(r,y,x)|^2>0\quad(x\ne0).
\]
Set
\begin{equation}\label{eq:normalizedchart}
 A(r,y)=\diag(\ell/\sqrt L),\qquad
 \tau(r,y,x)=T(r,y,x)/\sqrt{N(r,y,x)}.
\end{equation}
These formulas also apply at repeated or zero principal curvatures.

\subsection*{Polynomial identity}
Set $\mathscr D(r,y)=5625L^5/32>0$. For $x\ne0$, define
\begin{equation}\label{eq:Pidentity}
 P(r,y,x)=\mathscr D(r,y)N(r,y,x)^2
 \left[-W_{1/2}\bigl(A(r,y),\tau(r,y,x)\bigr)-\frac1{10000}\right].
\end{equation}
Substitution of \eqref{eq:endpointexpanded} cancels the denominators
and extends $P$ to an integer polynomial, homogeneous of degree four
in $x$ and of degree at most $14$ in each of $r,y$. It has 196 quartic
monomials and 19\,020 nonzero bivariate coefficient terms.

\subsection*{Bernstein coefficients and positivity}
Fix $r_0<r_1$ and $y_0<y_1$. On the rectangle
$[r_0,r_1]\times[y_0,y_1]$, set
\[
 s_1=\frac{r-r_0}{r_1-r_0},\qquad
 s_2=\frac{y-y_0}{y_1-y_0}.
\]
For $0\le i,j\le14$, define
\begin{equation}\label{eq:Bernsteinbasis}
 \beta_{ij}(r,y)=\binom{14}{i}s_1^i(1-s_1)^{14-i}
                 \binom{14}{j}s_2^j(1-s_2)^{14-j}.
\end{equation}
These nonnegative polynomials satisfy $\sum_{i,j=0}^{14}\beta_{ij}=1$.
Define the quartics $P_{ij}$ and scalars $d_{ij}$ by
\begin{equation}\label{eq:Bernstein}
 P=\sum_{i,j=0}^{14}\beta_{ij}P_{ij}(x),\qquad
 \mathscr D=\sum_{i,j=0}^{14}\beta_{ij}d_{ij}.
\end{equation}
Label each quadratic monomial $x_ax_b$ by $(a,b)$, with
$1\le a\le b\le12$, and let $z(x)\in\R^{78}$ list these monomials
in lexicographic order. It contains every $x_a^2$, so $z(x)\ne0$
whenever $x\ne0$.

\begin{lemma}[Bernstein positivity]\label{lem:certificate}
Fix a rectangle and the expansions \eqref{eq:Bernstein}.
Suppose that $d_{ij}>0$ for every $0\le i,j\le14$, and that for
each such pair there is a positive definite symmetric matrix
$G_{ij}\in\R^{78\times78}$ satisfying
\[
 P_{ij}(x)=d_{ij}\,z(x)^TG_{ij}z(x)\qquad(x\in\R^{12}).
\]
Then $P(r,y,x)>0$ for every point $(r,y)$ of the rectangle and
every $x\ne0$.
\end{lemma}
\begin{proof}
For $x\ne0$, each $P_{ij}(x)$ is positive. Equation~\eqref{eq:Bernstein}
is a weighted sum of these values with nonnegative weights of sum one.
\end{proof}

\begin{lemma}[Integer positivity test]\label{lem:integer}
Let $s\ge1$, let $Q\in\mathbb Z^{s\times s}$ be symmetric,
and let $V\in\mathbb Z^{s\times s}$. If
\begin{equation}\label{eq:integercriterion}
 (VQV^T)_{ii}>\sum_{j\ne i}|(VQV^T)_{ij}|\qquad(1\le i\le s),
\end{equation}
then $Q$ is positive definite.
\end{lemma}
\begin{proof}
Put $M=VQV^T$. For $w\in\R^s$, the inequality
$2|w_iw_j|\le w_i^2+w_j^2$ gives
\[
 w^TMw\ge\sum_i\left(M_{ii}-\sum_{j\ne i}|M_{ij}|\right)w_i^2.
\]
Hence $M$ is positive definite. If $V$ were singular, a nonzero
$w\in\ker V^T$ would satisfy $w^TMw=0$. Thus $V$ is invertible
and $Q$ is positive definite by congruence.
\end{proof}

\subsection*{The rational certificate}
A binary word selects closed halves of $[0,1]^2$: bisect in $r$ at
even depth and $y$ at odd depth, starting at depth zero, and retain
the lower half for 0 and the upper half for 1. The 24 leaves listed
in \texttt{DATA\_SPEC.md} form a complete prefix tree: every non-leaf
prefix has both children, and no leaf has a listed descendant.
Their rectangles therefore cover the closed square with disjoint interiors.

To prove $d_{ij}>0$, use $\mathscr D=180000(1+r^2+r^2y^2)^5$.
On any listed rectangle,
write $r=r_0+hs_1$, $y=y_0+ks_2$, where $r_0,y_0\ge0$ and $h,k>0$.
If $\mathscr D=\sum_{a,b}e_{ab}s_1^as_2^b$, then $e_{ab}\ge0$ and
$e_{00}>0$. The degree is at most ten in each variable. Conversion
to degree $(14,14)$ gives
\[
 d_{ij}=\sum_{a\le i,\,b\le j}e_{ab}
 \frac{\binom ia}{\binom{14}a}\frac{\binom jb}{\binom{14}b}>0.
\]
The positive term $e_{00}$ contributes to every coefficient.

Reversing a basis vector $e_k$ changes the sign of a tensor component
once for each occurrence of $k$ among its indices. Group the quadratic
monomials according to their signs under these four reversals.
This gives eight groups of sizes $16,9,9,9,9,9,9,8$. We take
$G_{ij}$ block diagonal with respect to these groups, each block
indexed by the monomials in that group. The encoding and ordering
are specified in \texttt{DATA\_SPEC.md}.
Index the 442 upper-triangular positions by $e$,
with weight $w_e=1$ on the diagonal and $w_e=2$ otherwise.
The sorted concatenation of the two quadratic labels gives a quartic
label $\gamma(e)$; there are 196 such labels. Set
\[
 c_\gamma=\sum_{\gamma(e)=\gamma}w_e.
\]
These labels give precisely the degree-four monomials unchanged by
all four sign reversals.
The other 1\,169 of the $\binom{15}{4}=1365$ quartic types have zero
coefficient; all are included in the coefficient checks.

At each rectangle and control index $(i,j)$, choose a common positive
integer scale $\sigma$ clearing the denominators of $P_{ij}$ and $d_{ij}$,
as specified in \texttt{DATA\_SPEC.md}. Let $p_\gamma$ be the coefficient
of $x^\gamma$ in $\sigma P_{ij}$, and let $d$ be the positive integer
$\sigma d_{ij}$.
Here $x^\gamma=x_{a_1}x_{a_2}x_{a_3}x_{a_4}$ if
$\gamma=(a_1,a_2,a_3,a_4)$. The canonical Gram entry is
$p_{\gamma(e)}/(c_{\gamma(e)}d)$. Each rectangle supplies
$n^{00},n^{10},n^{01},n^{11}\in\mathbb Z^{442}$ and the positive
denominator $n_*=12884901888$, with
\[
 \sum_{\gamma(e)=\gamma}w_e n_e^{ab}=0
 \quad\text{for every label $\gamma$ and }ab\in\{00,10,01,11\}.
\]
Define the integer vector
\[
 \nu(i,j)=(14-i)(14-j)n^{00}+i(14-j)n^{10}
                +(14-i)j n^{01}+ij n^{11}.
\]
The Gram entry at position $e$ is
\begin{equation}\label{eq:compressedgram}
 (G_{ij})_e=\frac{p_{\gamma(e)}}{c_{\gamma(e)}d}
                     +\frac{\nu_e(i,j)}{196n_*}.
\end{equation}
The first term supplies the coefficients of $P_{ij}/d_{ij}$;
the second represents the zero quartic. Thus
\begin{equation}\label{eq:Gramcertificate}
 P_{ij}(x)/d_{ij}=z(x)^TG_{ij}z(x).
\end{equation}
Positive definiteness is checked at every control index, not inferred
from the four corners.

The integer numerator $Q_{ij}=(196n_*d)G_{ij}$ has entries
\[
 (Q_{ij})_e=\frac{196n_*}{c_{\gamma(e)}}p_{\gamma(e)}
                      +d\,\nu_e(i,j).
\]
The possible multiplicities $c_\gamma$ are $1,2,3,4,6$, according
to the repeated indices in $\gamma$. Each divides $196n_*$, so
these entries are integers.
The positive factor $196n_*d$ preserves definiteness. For each block,
the data supply an integer matrix $V$ satisfying
\eqref{eq:integercriterion}. The Bernstein coefficients are computed
both by affine substitution and by midpoint restriction.
For each computation, the verification includes
$24\cdot15^2=5400$ full Gram identities, $43\,200$ block tests and
$5400\cdot78=421\,200$ strict row inequalities.

The verification consists of checking:
\begin{enumerate}
\item the full coefficient identity \eqref{eq:Pidentity}, reconstructed
both from \eqref{eq:K}--\eqref{eq:ZH} and from contractions on
$\Sym^4\R^4$ with the ordered inner product;
\item the 24-rectangle cover, the Bernstein expansions and
\eqref{eq:Gramcertificate} coefficient by coefficient, including
positivity of all denominators and all $d_{ij}$;
\item symmetry and the strict inequalities \eqref{eq:integercriterion}
after the supplied congruence, for all $43\,200$ integer blocks.
\end{enumerate}
These checks have been completed using exact rational arithmetic.
Lemma~\ref{lem:certificate} and
\eqref{eq:Pidentity} then prove Proposition~\ref{prop:endpoint}.

\subsection*{Electronic supplement}
The TeX source, Mathematica notebook and PDF export, and verification
code and data are at
\url{https://github.com/Jackyluckin/proof-of-second-pinching-prob-dim4}.
See \texttt{README.md} for instructions and
\texttt{DATA\_SPEC.md} for the data format.

The Lean version of Theorem~\ref{thm:main} is conditional on the finite
coefficient identities and integer positivity checks in \texttt{lean/README.md}.
Mathematica verifies them; their Lean proof is not included.
The classical classification results and Corollary~\ref{cor:main}
are not formalised.

\subsection*{Use of AI}
AI tools, including GPT-6 Astra, were used to explore candidate estimates,
develop the verification code, review the arguments, and revise the
manuscript. The author is responsible for the mathematical content and
its presentation.

\begin{thebibliography}{99}
\bibitem{Cartan}
\'{E}. Cartan,
\emph{Sur des familles remarquables d'hypersurfaces isoparam\'etriques dans les espaces sph\'eriques},
Math. Z. \textbf{45} (1939), 335--367,
\href{https://doi.org/10.1007/BF01580289}{\nolinkurl{doi:10.1007/BF01580289}}.

\bibitem{CDK}
S.-S. Chern, M. do Carmo and S. Kobayashi,
\emph{Minimal submanifolds of a sphere with second fundamental form of constant length},
in \emph{Functional Analysis and Related Fields}, F. E. Browder (ed.),
Springer, Berlin, Heidelberg, 1970, pp.~59--75,
\href{https://doi.org/10.1007/978-3-642-48272-4_2}{\nolinkurl{doi:10.1007/978-3-642-48272-4_2}}.

\bibitem{Chen}
R. Chen,
\emph{An improved second pinching theorem for minimal hypersurfaces with constant scalar curvature in spheres},
preprint (2026), \href{https://arxiv.org/abs/2609.17586v1}{arXiv:2609.17586v1}.

\bibitem{CLW}
Q.-M. Cheng, F. Li and G. Wei,
\emph{An improved result on Chern conjecture},
preprint (2026), \href{https://arxiv.org/abs/2610.05276v1}{\nolinkurl{arXiv:2610.05276v1}}.

\bibitem{GLLY}
J. Ge, T. Liu, K. Luo and W. Yan,
\emph{Rigidity of closed minimal hypersurfaces in $\Sph^5$},
preprint (2026), \href{https://arxiv.org/abs/2606.29246v1}{arXiv:2606.29246v1}.

\bibitem{Guan}
S. Guan,
\emph{The second gap for minimal hypersurfaces in the unit sphere with constant cubic trace},
preprint (2026), \href{https://arxiv.org/abs/2609.26852v1}{arXiv:2609.26852v1}.

\bibitem{GTYZ}
J. Ge, H. Tan, W. Yan and Y. Zhang,
\emph{The second gap and rigidity in Chern's conjecture with constant cubic trace},
preprint (2026), \href{https://arxiv.org/abs/2609.27661v1}{arXiv:2609.27661v1}.

\bibitem{HXZ}
C. He, H. Xu and E. Zhao,
\emph{Classification of closed minimal hypersurfaces with constant scalar curvature in $\Sph^5$},
preprint (2026), \href{https://arxiv.org/abs/2603.01181v1}{arXiv:2603.01181v1}.

\bibitem{Lawson}
H. B. Lawson, Jr.,
\emph{Local rigidity theorems for minimal hypersurfaces},
Ann. of Math. (2) \textbf{89} (1969), 187--197,
\href{https://doi.org/10.2307/1970816}{\nolinkurl{doi:10.2307/1970816}}.

\bibitem{LXX}
L. Lei, H. Xu and Z. Xu,
\emph{On Chern's conjecture for minimal hypersurfaces in spheres},
preprint (2017), \href{https://arxiv.org/abs/1712.01175v1}{arXiv:1712.01175v1}.

\bibitem{Munzner1}
H. F. M\"unzner,
\emph{Isoparametrische Hyperfl\"achen in Sph\"aren},
Math. Ann. \textbf{251} (1980), 57--71,
\href{https://doi.org/10.1007/BF01420281}{\nolinkurl{doi:10.1007/BF01420281}}.

\bibitem{Munzner2}
H. F. M\"unzner,
\emph{Isoparametrische Hyperfl\"achen in Sph\"aren. II.
\"Uber die Zerlegung der Sph\"are in Ballb\"undel},
Math. Ann. \textbf{256} (1981), 215--232,
\href{https://doi.org/10.1007/BF01450799}{\nolinkurl{doi:10.1007/BF01450799}}.

\bibitem{PT1}
C. K. Peng and C. L. Terng,
\emph{Minimal hypersurfaces of spheres with constant scalar curvature},
in \emph{Seminar on Minimal Submanifolds}, E. Bombieri (ed.),
Ann. of Math. Stud. \textbf{103},
Princeton University Press, Princeton, NJ, 1983, pp.~177--198.

\bibitem{PT2}
C. K. Peng and C. L. Terng,
\emph{The scalar curvature of minimal hypersurfaces in spheres},
Math. Ann. \textbf{266} (1983), 105--113,
\href{https://doi.org/10.1007/BF01458707}{\nolinkurl{doi:10.1007/BF01458707}}.

\bibitem{Simons}
J. Simons,
\emph{Minimal varieties in Riemannian manifolds},
Ann. of Math. (2) \textbf{88} (1968), 62--105,
\href{https://doi.org/10.2307/1970556}{\nolinkurl{doi:10.2307/1970556}}.

\bibitem{SY}
Y. J. Suh and H. Y. Yang,
\emph{The scalar curvature of minimal hypersurfaces in a unit sphere},
Commun. Contemp. Math. \textbf{9} (2007), no.~2, 183--200,
\href{https://doi.org/10.1142/S021919970700237X}{\nolinkurl{doi:10.1142/S021919970700237X}}.
\end{thebibliography}
\end{document}
